% HD-R12. GPT-6 Astra Pro /root; inherited effort not exposed.
% 21 November 2026. Exact coefficient-field and spectral-projector restriction.
\subsection{HD-R12: cyclotomic lower bounds for the actual coefficient rank}
\label{sec:cyclotomic-coefficient-rank}

For
\[
 p\in\{7,11,19,43,67,163\},\qquad K=\mathbb Q(\sqrt{-p}),
                                                               \tag{R12.1}
\]
every finite subgroup \(G\subset GL_m(K)\) whose character
field is exactly \(K\) satisfies
\[
                    m\ge\frac{p-1}{2}.                     \tag{R12.2}
\]
Absolute irreducibility is unnecessary for this bound.
For every actual element \(g\in G\) with
\(\operatorname{tr}(g)\notin\mathbb Q\), the stronger local
statement is
\[
               \operatorname{rank}_K(g-I)\ge\frac{p-1}{2}.
                                                               \tag{R12.3}
\]
The proof specifies the cyclotomic field inclusion, the Galois
restriction map, and an actual polynomial projector on the
coefficient space. No comparison of Jacobian factors is needed
for this restriction on their already specified coefficients.

\paragraph{1. The exact field extension detected by a characteristic factor.}
Let \(g\) have finite order \(h\). Its minimal polynomial
divides \(X^h-1\), which has distinct roots in characteristic
zero. The coefficient space is therefore a direct sum of its
actual irreducible \(K[X]\)-summands. All roots are roots of
unity. A nonrational trace is a nonrational coefficient of
\(\det(XI-g)\), so at least one of its monic irreducible
factors over \(K\), say \(f\), has nonrational coefficients.
Choose a root \(\zeta\) of \(f\), of exact order \(N\),
and put \(F=\mathbb Q(\zeta)\), \(L=KF\).

The actual homomorphism
\[
 \operatorname{Gal}(L/K)\longrightarrow
       \operatorname{Gal}(F/(K\cap F)),\qquad
                   \sigma\longmapsto\sigma|_F              \tag{R12.4}
\]
is an isomorphism. One direct verification in this quadratic
situation is as follows. If \(K\subset F\), then \(L=F\)
and the map is the identity on the indicated subgroup. Otherwise
the polynomial \(X^2+p\) remains irreducible over \(F\),
so \([L:F]=2\) and \([L:K]=[F:\mathbb Q]\). The restriction
map is injective because \(K\) and \(F\) generate \(L\),
and both sides then have that same order. This proves the claim
and explicitly describes its inverse as the unique extension
fixing \(K\).

In the second case the orbit of \(\zeta\) is its entire
rational cyclotomic orbit. The minimal polynomial \(f\) is
then \(\Phi_N\in\mathbb Q[X]\), contrary to its choice.
We have therefore obtained the literal inclusion
\[
                     K\subset\mathbb Q(\zeta_N),\qquad
                  \deg f=[F:K]=\frac{\varphi(N)}2.          \tag{R12.5}
\]
Here the cyclotomic degree and irreducibility are the published
Theorem6.4(a) in J.~S.~Milne, \emph{Algebraic Number Theory},
version3.08 (19 July2020), printed p.99. Its part(c), on the
same page, gives the exact ramification statement used next.
Primary source: \url{https://www.jmilne.org/math/CourseNotes/ANT.pdf}.

\paragraph{2. Ramification and the actual moving summand.}
Write \(s=\sqrt{-p}\). The integral principal ideal \((s)\)
of \(\mathcal O_K\) has norm \(p\), hence is prime, and
\((s)^2=(p)\). Thus \(p\) ramifies in \(K\). The inclusion
in (R12.5), followed by the inclusion into the same complex
algebraic closure, is a tower of actual number fields. Extension
of valuations multiplies its ramification indices. Therefore
\(p\) ramifies in \(F\). Milne's Theorem6.4(c) gives
\(p\mid N\). Writing \(N=p^aM\), \((p,M)=1\), yields
\[
         \varphi(N)=p^{a-1}(p-1)\varphi(M)\ge p-1.         \tag{R12.6}
\]

To retain the coefficient subspace explicitly, let \(\mu\)
be the square-free minimal polynomial of \(g\), and put
\(q=\mu/f\). Euclidean division in \(K[X]\) gives
\(u q+v f=1\). The actual endomorphism
\[
                     e_f=u(g)q(g)                         \tag{R12.7}
\]
is the identity on \(\ker f(g)\), zero on every other
irreducible summand, and satisfies \(e_f^2=e_f\).
The nonzero vector space \(e_fV\) is a module over the
explicit field \(K[X]/(f)\), via \(X\mapsto g|_{e_fV}\).
For a basis \(v_1,\ldots,v_a\) over that field, the map
\[
 (K[X]/(f))^a\longrightarrow e_fV,\qquad
      (r_1,\ldots,r_a)\longmapsto\sum_j r_j(g)v_j           \tag{R12.8}
\]
is a \(K\)-linear isomorphism intertwining multiplication
by \(X\) and the original matrix \(g\).
Since \(f\) has nonrational coefficients, it is not
\(X-1\). Multiplication by \(X-1\) on this field is
invertible. Consequently
\[
 \operatorname{rank}_K(g-I)\ge\dim_K(e_fV)
            =a\deg f\ge\frac{\varphi(N)}2\ge\frac{p-1}{2}.
                                                               \tag{R12.9}
\]
This proves (R12.3). Exact character field \(K\) means that
some element has nonrational trace; applying (R12.9) to it
proves (R12.2). The element need not be one of the selected
inertias, and no assertion of that kind is used.

\paragraph{3. The resulting main-source gate.}
Combining this coefficient restriction with the complete
rank-at-most-fourteen source closure in HD-S40 and HD-G28
gives the following necessary ranks for sources of the missing
sixfold determinant components:
\[
\begin{array}{c|rrrrrrrr}
 d&1&2&7&11&19&43&67&163\\\hline
 \text{remaining coefficient rank at least}
   &15&15&15&15&15&21&33&81.
\end{array}                                                   \tag{R12.10}
\]
For \(d=1\) the missing determinant classes are outside
\([-1],[-3]\); for the other seven fields they are outside
\([-1]\). The Eisenstein family already has its separate
effective supplier construction. Every remaining source here is
reflection-free by HD-S36 and HD-G27. Formula (R12.10) is a
necessary coefficient restriction, without an existence or
dominance claim at any listed endpoint.

The field \(K[X]/(f)\) in (R12.8) is also a finite etale
\(K\)-algebra. Its multiplication module describes the spectral
summand after the literal restriction of the given representation
to \(H=\langle g\rangle\). In the curve-source setting, write
\(\widetilde U\to U\) for the actual finite etale full-\(G\)
cover of the punctured base. The finite etale intermediate map
\[
                   U_H=\widetilde U/H\longrightarrow U
                                                               \tag{R12.11}
\]
pulls the coefficient sheaf back to this restricted representation,
after fixing the usual basepoint and sheet. The projector \(e_f\)
commutes with \(H\), so it is an actual idempotent endomorphism
of that pulled-back sheaf and its image is the asserted spectral
sub-local-system there. No \(G\)-invariance of \(e_fV\) on
the original base is asserted. Its characteristic factor, exterior
powers and traces are preserved under each coefficient embedding
into an algebraic closure of \(\mathbb Q_\ell\). These are
the precise coefficient and finite-cover maps related to the
realizations described in the framework section.
It restricts the curve suppliers for condition A (cote16,
scan050, with SIGN\_NEXT SN-C03/C04); it introduces no effective
cycle to which the existing condition-I and resultant-projector
division steps could yet be applied on a missing component.
